%! TEX root = thesis.tex

\chapter{Validation and Verification}
\label{sec:validation_verification}

\section{Validation}

Validation checks whether the simulation is an adequate representation of
real-world carbon-aware training for the study objectives.  It emphasizes
consistency with domain expectations rather than low-level code correctness.

\subsection{Structural Trade-off Behaviour}

The first validation checks the core trade-off between CO$_2$ savings and
training time overhead when varying the pause threshold $\theta_p$.  For fixed
region, start time, and model configuration, stricter thresholds (lower
$\theta_p$) should yield higher relative emission savings but also higher
wall-clock overhead.  The simulator is validated by confirming a monotone
trend in this trade-off and by inspecting the Pareto frontier for reasonable
shape and absence of pathological cases, such as negative savings for very
strict thresholds.


\subsection{Regional and Temporal Realism}

The third validation targets regional and temporal realism.  For each region
and year, an average grid carbon intensity is derived from the input time
series, and baseline emissions in a no-pause scenario should broadly follow
this ordering for fixed IT energy.  Additionally, under a chosen carbon-aware
pausing policy, regions or periods with high temporal variability in CO$_2$
intensity should exhibit greater potential savings than low-variance,
already-clean regions, albeit with more complex pause patterns.  Anomalous
results are checked against features of the underlying time series to ensure
they can be explained by specific training windows.


\section{Verification}

Verification ensures that the implementation follows the intended mathematical
model and uses the specified inputs and state updates correctly.  It is
carried out via targeted unit tests and synthetic test cases.

\subsection{Energy and Emissions Accounting}

The first verification step checks the energy and emissions accounting.  At
each time step, IT power, data center power, and incremental emissions are
computed from GPU counts, training or pause power, power usage effectiveness, and the CO$_2$
intensity.  The accumulated emissions over time, as well as baseline and
policy totals, are compared against analytical results in simple synthetic
scenarios.  Additionally, the KPI formulas for CO$_2$ savings percentage, time
overhead percentage, and the composite score are verified to match their
definitions for fixed numerical inputs.

\subsection{Pause/Resume Logic and Hysteresis}

The second verification step focuses on the pause/resume logic driven by the
pause threshold $\theta_p$ and hysteresis $\theta_r$.  A synthetic sequence of
CO$_2$ intensities is constructed to trigger known patterns of pausing and
resuming according to the specified hysteresis rule.  The simulated pause
state and pause count at each step are compared to the expected sequence,
ensuring that transitions only occur when the intensity crosses the configured
thresholds in the intended direction.

\subsection{Progress and Runtime Consistency}

The third verification step addresses consistency between training progress
and wall-clock time.  In non-paused steps, the progress fraction should
increase in proportion to the time step relative to the baseline training
time, while the elapsed wall time grows in every step.  In a no-pause
scenario, the simulation is expected to reach full progress exactly at the
baseline runtime within numerical tolerance.  In scenarios with pauses, the
simulated policy runtime is compared against the analytical expression
combining baseline time, total paused time, and per-pause checkpoint overhead.

\subsection{Unit Tests for KPI Functions}

Finally, KPI computations are tested in isolation through unit tests.  Fixed
values are supplied for baseline and policy emissions and runtimes, and the
functions computing CO$_2$ savings percentage, time overhead percentage, and
the composite score are checked against manually computed expected results.
This confirms that the KPI implementations are numerically correct,
independent of the rest of the simulation engine.

